% Options for packages loaded elsewhere
\PassOptionsToPackage{unicode}{hyperref}
\PassOptionsToPackage{hyphens}{url}
\documentclass[
]{article}
\usepackage{xcolor}
\usepackage[margin=1in]{geometry}
\usepackage{amsmath,amssymb}
\setcounter{secnumdepth}{-\maxdimen} % remove section numbering
\usepackage{iftex}
\ifPDFTeX
  \usepackage[T1]{fontenc}
  \usepackage[utf8]{inputenc}
  \usepackage{textcomp} % provide euro and other symbols
\else % if luatex or xetex
  \usepackage{unicode-math} % this also loads fontspec
  \defaultfontfeatures{Scale=MatchLowercase}
  \defaultfontfeatures[\rmfamily]{Ligatures=TeX,Scale=1}
\fi
\usepackage{lmodern}
\ifPDFTeX\else
  % xetex/luatex font selection
\fi
% Use upquote if available, for straight quotes in verbatim environments
\IfFileExists{upquote.sty}{\usepackage{upquote}}{}
\IfFileExists{microtype.sty}{% use microtype if available
  \usepackage[]{microtype}
  \UseMicrotypeSet[protrusion]{basicmath} % disable protrusion for tt fonts
}{}
\makeatletter
\@ifundefined{KOMAClassName}{% if non-KOMA class
  \IfFileExists{parskip.sty}{%
    \usepackage{parskip}
  }{% else
    \setlength{\parindent}{0pt}
    \setlength{\parskip}{6pt plus 2pt minus 1pt}}
}{% if KOMA class
  \KOMAoptions{parskip=half}}
\makeatother
\usepackage{color}
\usepackage{fancyvrb}
\newcommand{\VerbBar}{|}
\newcommand{\VERB}{\Verb[commandchars=\\\{\}]}
\DefineVerbatimEnvironment{Highlighting}{Verbatim}{commandchars=\\\{\}}
% Add ',fontsize=\small' for more characters per line
\usepackage{framed}
\definecolor{shadecolor}{RGB}{248,248,248}
\newenvironment{Shaded}{\begin{snugshade}}{\end{snugshade}}
\newcommand{\AlertTok}[1]{\textcolor[rgb]{0.94,0.16,0.16}{#1}}
\newcommand{\AnnotationTok}[1]{\textcolor[rgb]{0.56,0.35,0.01}{\textbf{\textit{#1}}}}
\newcommand{\AttributeTok}[1]{\textcolor[rgb]{0.13,0.29,0.53}{#1}}
\newcommand{\BaseNTok}[1]{\textcolor[rgb]{0.00,0.00,0.81}{#1}}
\newcommand{\BuiltInTok}[1]{#1}
\newcommand{\CharTok}[1]{\textcolor[rgb]{0.31,0.60,0.02}{#1}}
\newcommand{\CommentTok}[1]{\textcolor[rgb]{0.56,0.35,0.01}{\textit{#1}}}
\newcommand{\CommentVarTok}[1]{\textcolor[rgb]{0.56,0.35,0.01}{\textbf{\textit{#1}}}}
\newcommand{\ConstantTok}[1]{\textcolor[rgb]{0.56,0.35,0.01}{#1}}
\newcommand{\ControlFlowTok}[1]{\textcolor[rgb]{0.13,0.29,0.53}{\textbf{#1}}}
\newcommand{\DataTypeTok}[1]{\textcolor[rgb]{0.13,0.29,0.53}{#1}}
\newcommand{\DecValTok}[1]{\textcolor[rgb]{0.00,0.00,0.81}{#1}}
\newcommand{\DocumentationTok}[1]{\textcolor[rgb]{0.56,0.35,0.01}{\textbf{\textit{#1}}}}
\newcommand{\ErrorTok}[1]{\textcolor[rgb]{0.64,0.00,0.00}{\textbf{#1}}}
\newcommand{\ExtensionTok}[1]{#1}
\newcommand{\FloatTok}[1]{\textcolor[rgb]{0.00,0.00,0.81}{#1}}
\newcommand{\FunctionTok}[1]{\textcolor[rgb]{0.13,0.29,0.53}{\textbf{#1}}}
\newcommand{\ImportTok}[1]{#1}
\newcommand{\InformationTok}[1]{\textcolor[rgb]{0.56,0.35,0.01}{\textbf{\textit{#1}}}}
\newcommand{\KeywordTok}[1]{\textcolor[rgb]{0.13,0.29,0.53}{\textbf{#1}}}
\newcommand{\NormalTok}[1]{#1}
\newcommand{\OperatorTok}[1]{\textcolor[rgb]{0.81,0.36,0.00}{\textbf{#1}}}
\newcommand{\OtherTok}[1]{\textcolor[rgb]{0.56,0.35,0.01}{#1}}
\newcommand{\PreprocessorTok}[1]{\textcolor[rgb]{0.56,0.35,0.01}{\textit{#1}}}
\newcommand{\RegionMarkerTok}[1]{#1}
\newcommand{\SpecialCharTok}[1]{\textcolor[rgb]{0.81,0.36,0.00}{\textbf{#1}}}
\newcommand{\SpecialStringTok}[1]{\textcolor[rgb]{0.31,0.60,0.02}{#1}}
\newcommand{\StringTok}[1]{\textcolor[rgb]{0.31,0.60,0.02}{#1}}
\newcommand{\VariableTok}[1]{\textcolor[rgb]{0.00,0.00,0.00}{#1}}
\newcommand{\VerbatimStringTok}[1]{\textcolor[rgb]{0.31,0.60,0.02}{#1}}
\newcommand{\WarningTok}[1]{\textcolor[rgb]{0.56,0.35,0.01}{\textbf{\textit{#1}}}}
\usepackage{longtable,booktabs,array}
\newcounter{none} % for unnumbered tables
\usepackage{calc} % for calculating minipage widths
% Correct order of tables after \paragraph or \subparagraph
\usepackage{etoolbox}
\makeatletter
\patchcmd\longtable{\par}{\if@noskipsec\mbox{}\fi\par}{}{}
\makeatother
% Allow footnotes in longtable head/foot
\IfFileExists{footnotehyper.sty}{\usepackage{footnotehyper}}{\usepackage{footnote}}
\makesavenoteenv{longtable}
\usepackage{graphicx}
\makeatletter
\newsavebox\pandoc@box
\newcommand*\pandocbounded[1]{% scales image to fit in text height/width
  \sbox\pandoc@box{#1}%
  \Gscale@div\@tempa{\textheight}{\dimexpr\ht\pandoc@box+\dp\pandoc@box\relax}%
  \Gscale@div\@tempb{\linewidth}{\wd\pandoc@box}%
  \ifdim\@tempb\p@<\@tempa\p@\let\@tempa\@tempb\fi% select the smaller of both
  \ifdim\@tempa\p@<\p@\scalebox{\@tempa}{\usebox\pandoc@box}%
  \else\usebox{\pandoc@box}%
  \fi%
}
% Set default figure placement to htbp
\def\fps@figure{htbp}
\makeatother
\setlength{\emergencystretch}{3em} % prevent overfull lines
\providecommand{\tightlist}{%
  \setlength{\itemsep}{0pt}\setlength{\parskip}{0pt}}
\usepackage{bookmark}
\IfFileExists{xurl.sty}{\usepackage{xurl}}{} % add URL line breaks if available
\urlstyle{same}
\hypersetup{
  pdftitle={Testing out SQL},
  hidelinks,
  pdfcreator={LaTeX via pandoc}}

\title{Testing out SQL}
\author{}
\date{\vspace{-2.5em}2026-08-03}

\begin{document}
\maketitle

\begin{Shaded}
\begin{Highlighting}[]
\FunctionTok{library}\NormalTok{(DBI) }
\FunctionTok{library}\NormalTok{(RMariaDB)}
\end{Highlighting}
\end{Shaded}

\section{A beginners guide to SQL.}\label{a-beginners-guide-to-sql.}

Testing out different functions, using Rmarkdown

\subsection{Connecting SQL i
Rmarkdown}\label{connecting-sql-i-rmarkdown}

This runs the SQL-connection named con in R, Password is a hidden param.
Not relevant for SQL in itself but needed for connecting to the database
in Rmarkdown.

\begin{Shaded}
\begin{Highlighting}[]
\NormalTok{con }\OtherTok{\textless{}{-}} \FunctionTok{dbConnect}\NormalTok{(}
\NormalTok{  RMariaDB}\SpecialCharTok{::}\FunctionTok{MariaDB}\NormalTok{(),}
  \AttributeTok{dbname =} \StringTok{"test"}\NormalTok{,}
  \AttributeTok{host =} \StringTok{"localhost"}\NormalTok{,}
  \AttributeTok{user =} \StringTok{"root"}\NormalTok{,}
  \AttributeTok{password =}\NormalTok{ pw}
\NormalTok{)}
\end{Highlighting}
\end{Shaded}

\subsection{Create a table}\label{create-a-table}

Constraints: \textbf{NOT NULL}.the value needs to be not null,
\textbf{UNIQUE} needs to be unique value

\textbf{DEFAULT} sätter ett defaultvärde vid null.
\textbf{AUTO\_INCREMENT} gör att du inte behöver skriva in student\_id
den räknar själv 1,2,3 o.s.v. ON DELETE sätter beteendet vid
borttagandet om foreign keys används i andra tabeller. \textbf{ON DELETE
SET NULL} sätter värdet i de andra tabellerna till NULL. \textbf{ON
DELETE CASCADE} tar bort raden även i de andra tabellerna där foreign
key används. Exempel: FOREIGN KEY (major\_id) REFERENCES
major(major\_id) ON DELETE SET NULL

\begin{Shaded}
\begin{Highlighting}[]
\KeywordTok{CREATE} \KeywordTok{TABLE}\NormalTok{ student (}
\NormalTok{    student\_id }\DataTypeTok{INT} \KeywordTok{PRIMARY} \KeywordTok{KEY}\NormalTok{ AUTO\_INCREMENT,}
\NormalTok{    name }\DataTypeTok{VARCHAR}\NormalTok{(}\DecValTok{20}\NormalTok{) }\KeywordTok{NOT} \KeywordTok{NULL}\NormalTok{,}
\NormalTok{    major }\DataTypeTok{VARCHAR}\NormalTok{(}\DecValTok{20}\NormalTok{) }\KeywordTok{DEFAULT} \OtherTok{"undecided"}
\NormalTok{);}
\end{Highlighting}
\end{Shaded}

\subsection{Describes a table}\label{describes-a-table}

\begin{Shaded}
\begin{Highlighting}[]
\NormalTok{DESCRIBE student;}
\end{Highlighting}
\end{Shaded}

\begin{longtable}[]{@{}llllll@{}}
\caption{3 records}\tabularnewline
\toprule\noalign{}
Field & Type & Null & Key & Default & Extra \\
\midrule\noalign{}
\endfirsthead
\toprule\noalign{}
Field & Type & Null & Key & Default & Extra \\
\midrule\noalign{}
\endhead
\bottomrule\noalign{}
\endlastfoot
student\_id & int & NO & PRI & NA & auto\_increment \\
name & varchar(20) & NO & & NA & \\
major & varchar(20) & YES & & undecided & \\
\end{longtable}

\subsection{Add a column}\label{add-a-column}

Vi lägger till en kolumn average som kan ha värdet \textbf{DECIMAL} där
vi anger 3 siffror totalt, 2st efter kommatecknet. ex 3.14

\begin{Shaded}
\begin{Highlighting}[]
\KeywordTok{ALTER} \KeywordTok{TABLE}\NormalTok{ student }\KeywordTok{ADD}\NormalTok{ average }\DataTypeTok{DECIMAL}\NormalTok{(}\DecValTok{3}\NormalTok{,}\DecValTok{2}\NormalTok{);}
\end{Highlighting}
\end{Shaded}

\subsection{Removes a column}\label{removes-a-column}

\begin{Shaded}
\begin{Highlighting}[]
\KeywordTok{ALTER} \KeywordTok{TABLE}\NormalTok{ student }\KeywordTok{DROP} \KeywordTok{COLUMN}\NormalTok{ average;}
\end{Highlighting}
\end{Shaded}

\subsection{Insert row to table.}\label{insert-row-to-table.}

In order of columns in table. Primary key som valdes i create table
måste vara unik.

\begin{Shaded}
\begin{Highlighting}[]
\KeywordTok{INSERT} \KeywordTok{INTO}\NormalTok{ student }\KeywordTok{VALUES}\NormalTok{(}\DecValTok{1}\NormalTok{, }\OtherTok{"Jack"}\NormalTok{, }\OtherTok{"Sociology"}\NormalTok{);}
\end{Highlighting}
\end{Shaded}

\section{Insert row to table ver 2}\label{insert-row-to-table-ver-2}

Specificera vilka värden du vill skriva in, om något ska utelämnas.
Första värdet student\_id behöver inte skrivas ut då vi använde
AUTO\_INCREMENT tidigare vid create table. Vi utelämnar även major,
vilket då bestäms till undecided av DEFAULT vi satte tidigare i CREATE
TABLE.

\begin{Shaded}
\begin{Highlighting}[]
\KeywordTok{INSERT} \KeywordTok{INTO}\NormalTok{ student(name) }\KeywordTok{VALUES}\NormalTok{(}\OtherTok{"Kate"}\NormalTok{);}
\end{Highlighting}
\end{Shaded}

\subsection{SELECT information from
table}\label{select-information-from-table}

\begin{itemize}
\tightlist
\item
  Används för att söka på allt. Ersätts t.ex. med name för att returnera
  alla namn. \textbf{SELECT DISTINCT} ger unika värden. \textbf{ORDER
  BY} name DESC; kan läggas till för att sortera t.ex. namn alfabetiskt,
  desc för omvänd ordning. \textbf{LIMIT 2}; returnerar max 2 värden.
\end{itemize}

\begin{Shaded}
\begin{Highlighting}[]
\KeywordTok{SELECT} \OperatorTok{*} 
\KeywordTok{FROM}\NormalTok{ student}
\KeywordTok{WHERE}\NormalTok{ name }\KeywordTok{IN}\NormalTok{ (}\OtherTok{"Kate"}\NormalTok{, }\OtherTok{"Jack"}\NormalTok{);}
\end{Highlighting}
\end{Shaded}

\begin{longtable}[]{@{}lll@{}}
\caption{2 records}\tabularnewline
\toprule\noalign{}
student\_id & name & major \\
\midrule\noalign{}
\endfirsthead
\toprule\noalign{}
student\_id & name & major \\
\midrule\noalign{}
\endhead
\bottomrule\noalign{}
\endlastfoot
1 & Jack & Sociology \\
2 & Kate & undecided \\
\end{longtable}

\subsection{UPDATE values}\label{update-values}

Ändrar alla värden som har ex. Sociology eller social psychology som
major till Soc. Även SET kan ha flera argument, lägg till efter ett
kommatecken. \textbf{\textless\textgreater{}} betyder not equal to

\begin{Shaded}
\begin{Highlighting}[]
\KeywordTok{UPDATE}\NormalTok{ student}
\KeywordTok{SET}\NormalTok{ major }\OperatorTok{=} \OtherTok{"Soc"}
\KeywordTok{WHERE}\NormalTok{ major }\OperatorTok{=} \OtherTok{"Sociology"} \KeywordTok{OR}\NormalTok{ major }\OperatorTok{=} \OtherTok{"Social psychology"}\NormalTok{ ;}
\end{Highlighting}
\end{Shaded}

\subsection{Delete values}\label{delete-values}

\begin{Shaded}
\begin{Highlighting}[]
\KeywordTok{DELETE} \KeywordTok{FROM}\NormalTok{ student}
\KeywordTok{WHERE}\NormalTok{ student\_id }\OperatorTok{=} \DecValTok{1}\NormalTok{;}
\end{Highlighting}
\end{Shaded}

\subsection{Delete a table}\label{delete-a-table}

\begin{Shaded}
\begin{Highlighting}[]
\KeywordTok{DROP} \KeywordTok{TABLE}\NormalTok{ student;}
\end{Highlighting}
\end{Shaded}

\section{Foreign keys}\label{foreign-keys}

En kolumn som pekar på en rad i en annan tabell -- eller på en rad i
samma tabell (självrefererande f k) Antingen via FOREIGN KEY
(student\_id) REFERENCES student(student\_id) i \textbf{CREATE TABLE}
eller via \textbf{ALTER TABLE} ex: ALTER TABLE grades ADD CONSTRAINT
fk\_student FOREIGN KEY (student\_id) REFERENCES student(student\_id);

\section{Wildcards}\label{wildcards}

Fungerar som enkla regular expressions. används t.ex. \textbf{WHERE name
LIKE ``xxx'';} \textbf{LIKE} \% betyder fler tecken och \_ betyder ett
tecken. ex: \%school\% ordet ska ha school i sig. \_\_-10\% betyder 2st
valfria tecken sedan -10 och valfritt antal tecken efter.

\section{UNION}\label{union}

Slå ihop resultat från två SELECT‑frågor till en gemensam lista.Använd
\textbf{UNION} mellan SELECT-frågorna.

\section{JOIN}\label{join}

Koppla ihop rader från två tabeller baserat på ett gemensamt värde.
\textbf{JOIN, LEFT JOIN, RIGHT JOIN \& FULL OUTER JOIN} join visar bara
värden som finns i båda, left join visar i detta exempel allt från
employee, right join visar allt från branch: SELECT employee.emp\_id,
branch.branch\_name FROM employee JOIN branch ON employee.emp\_id =
branch.mg\_id;

\section{TRIGGERS}\label{triggers}

\subsection{Triggers är till för att automatiskt reagera på förändringar
i en
tabell.}\label{triggers-uxe4r-till-fuxf6r-att-automatiskt-reagera-puxe5-fuxf6ruxe4ndringar-i-en-tabell.}

\textbf{Automatisera loggning}, Exempel: Spara gamla värden i en
historiktabell när något uppdateras. \textbf{Validera data}, Kontrollera
att en rad uppfyller vissa regler innan den får sparas. \textbf{Hålla
data synkroniserad}, Uppdatera relaterade tabeller när något ändras.
\textbf{Beräkna eller sätta värden automatiskt}, Exempel: sätta
updated\_at vid varje UPDATE. \textbf{Förhindra otillåtna ändringar},
Stoppa DELETE eller UPDATE om villkor inte uppfylls.

\end{document}
